% TOP_PROBLEM: {"id": 2302046, "title": "Research Problems in Function Theory — Problem 2.46", "category_id": 8, "category": {"id": 8, "name": "analysis", "display_name": "Analysis", "description": "Limits, continuity, calculus, and function theory.", "slug": "analysis", "order_index": 8, "created_at": "2026-07-31T15:26:25.671Z"}, "status": "open"}
% FIRST_POSTED: null
% ENTRY_KIND: statement_only
% Imported from: batch11.tex; UnsolvedMath ID 2302046
% Compile twice: lualatex 2302046.tex
\documentclass[11pt]{article}
\usepackage{amsmath,amssymb,mathtools}
\usepackage{fontspec}
\usepackage{unicode-math}
\setmainfont{Latin Modern Roman}
\setmathfont{Latin Modern Math}
\usepackage{xurl}
\usepackage[unicode,colorlinks=true,linkcolor=blue,urlcolor=blue,pdftitle={UnsolvedMath Problem 2302046}]{hyperref}
\newcommand{\sref}[2]{\hyperlink{source:#1:#2}{[S#2]}}
\newcommand{\cataloguescope}[1]{\paragraph{Mathematical statement.}#1}
\begin{document}
% UnsolvedMath ID 2302046; source catalogue code AMR-022-2046
\setcounter{section}{2302045}
\section{Research Problems in Function Theory — Problem 2.46}\label{2302046}
{\small\sffamily UnsolvedMath ID 2302046 · Analysis}
\subsection{Definitions and mathematical statement}

\paragraph{Shared notation.}$\mathbb N=\{1,2,\ldots\}$, and $\mathbb Z,\mathbb Q,\mathbb R,\mathbb C$ denote the integers, rationals, reals and complex numbers. Any other conventions are specified locally.


Let $\mathfrak c=|\mathbb R|=2^{\aleph_0}$ be the continuum cardinal, let $m$ be an infinite cardinal with $\aleph_0<m<\mathfrak c$, and let $\mathcal A$ be a set of distinct entire functions. Write $m^+$ for the successor cardinal of $m$.

\paragraph{Statement recorded in the supplied source.}
If
\[\bigl|\{f(z):f\in\mathcal A\}\bigr|\le m\quad\text{for every }z\in\mathbb C,\]
must $|\mathcal A|\le m$? In particular, determine the case $\mathfrak c=m^+$, or whether its answer is independent of the usual axioms of set theory. \sref{2302046}{2}
The original paragraph gives the positive answer when $m^+<\mathfrak c$, and motivates the question with Erdős's countable-family result and its continuum-hypothesis counterexample. Update 2.46 reports no further progress. \sref{2302046}{2}
\subsection{Short English statement}
If an entire-function family takes at most a specified number of different values at each point, must the family itself have at most that cardinality?
\subsection{Sources}
\begin{description}
\item[\hypertarget{source:2302046:1}{S1}] ULAM.AI and the UnsolvedMath Contributors, \emph{UnsolvedMath: A Curated Collection of Open Mathematics Problems}, record 2302046 (AMR-022-2046). CC BY 4.0. \url{https://huggingface.co/datasets/ulamai/UnsolvedMath}.
\item[\hypertarget{source:2302046:2}{S2}] W. K. Hayman and E. F. Lingham, \emph{Research Problems in Function Theory} (2018), Chapter 2, Problem 2.46 and Update 2.46. \url{https://arxiv.org/abs/1809.07200}.
\end{description}
\label{end:2302046}
\end{document}
